\section{ALCANCES Y LÍMITES}

\subsection{ALCANCES}
\begin{itemize}
	\item El diseño de un protocolo de votación electrónica que incorpora las primitivas criptográficas post-cuánticas especificadas en los estándares FIPS 203 y FIPS 204.
	\item La fundamentación teórica del protocolo a partir de la dificultad de los problemas Module-LWE y Module-SIS.
	\item El modelado formal del protocolo mediante Tamarin Prover bajo el modelo de atacante Dolev-Yao.
	\item La verificación de las propiedades de confidencialidad, integridad y autenticidad mediante lemmas formales.
	\item La implementación del protocolo con fines de simulación, orientada a medir indicadores de rendimiento --- latencia, tiempo de ejecución y consumo de memoria --- bajo distintos números de votantes.
	\item La comparación de dichos indicadores y propiedades de seguridad con un protocolo equivalente basado en RSA/ECDSA.
	\item El alcance metodológico se limita al análisis de la capa criptográfica del proceso de votación electrónica.
\end{itemize}

\subsection{LÍMITES}
\begin{itemize}
	\item El análisis formal y simbólico del protocolo mediante Tamarin Prover, en el cual las operaciones criptográficas especificadas en los estándares FIPS 203 y FIPS 204 se abstraen simbólicamente, sin modelar la estructura matemática subyacente de los retículos algebraicos.
	\item La implementación del protocolo no constituye un sistema de votación electrónica desplegado en un proceso electoral real, sino un entorno de simulación orientado a la medición de indicadores de rendimiento.
	\item La arquitectura de almacenamiento o registro distribuido de los votos no forma parte del estudio.
	\item La reproducibilidad exacta de las mediciones de rendimiento no está garantizada, dado que factores como el conjunto de instrucciones del procesador y las optimizaciones de la biblioteca afectan significativamente la latencia de los algoritmos post-cuánticos (Kannwischer et al., 2019).
	\item La comparación con el protocolo clásico se restringe a un esquema basado en RSA/ECDSA, sin extenderse a otros esquemas criptográficos clásicos o post-cuánticos adicionales.
\end{itemize}

